\documentclass[10pt,a4paper]{amsart}
\usepackage[margin=19mm]{geometry}
\usepackage{amsmath,amssymb,mathtools}
\usepackage[scr=rsfs,cal=euler]{mathalfa}
\usepackage{xfrac}
\usepackage[unicode,hidelinks,bookmarks=false]{hyperref}
\newcommand{\ie}{i.e.\ }
\newcommand{\cf}{cf.\ }
\newcommand{\resp}{respectively\ }
\newcommand{\Subsection}{\subsection}
\providecommand{\nobreakdash}{\nobreak\hbox{-}}
\setlength{\emergencystretch}{2em}
\multlinegap0pt
\usepackage{float}
\usepackage{amssymb}
\usepackage{braket}
\usepackage{mathtools}
\usepackage{tikz}
\newtheorem{lemma}{Lemma}
\newcommand{\e}{\operatorname{e}}
\title{Exponential sums over M\"{o}bius convolutions with applications to partitions}
\author{Debmalya Basak, Nicolas Robles, Alexandru Zaharescu}
\date{Source: arXiv:2312.17435v2 (CC BY 4.0)}
\begin{document}
\maketitle

\noindent\emph{Source and licence.} Exponential sums over M\"{o}bius convolutions with applications to partitions. Version arXiv:2312.17435v2 (CC BY 4.0), distributed by arXiv under the Creative Commons Attribution 4.0 International licence, http://creativecommons.org/licenses/by/4.0/, as recorded in the arXiv licence metadata for this exact version and shown on the abstract page https://arxiv.org/abs/2312.17435v2. Full licence notice: \texttt{licenses/arxiv-2312.17435-CC-BY-4.0.txt}.\par\smallskip

\noindent\textit{Editorial scope.} Counted unit: the article's lemma Theorem (source0.tex lines 956--961). The article's complete original proof of it is delivered with it (source0.tex lines 962--1005, one \texttt{proof} environment of 44 lines). No mathematics is written, repaired or re-derived: the statement and the proof are the article's own lines, in the article's own notation, and no retained line is substituted or re-flowed.  Retained uncounted as the source-local context the proof needs: lines 285--287; lines 860--862; lines 930--935.  Nothing was omitted from any retained span, so the package carries no omission record.  No retained line contains a citation, so the wrapper carries no bibliography and no bibliography entry.  No retained line contains a figure environment, an image-inclusion command or drawing code, so no image is delivered and the package declares no asset dependency.  Editorial formatting only, in addition: an AMS \texttt{amsart} preamble that restates the article's own declarations and macros used by the retained lines, namely the environments \texttt{lemma} on separate counters, together with the standard packages those lines need.  The retained environments print in this document as Lemma 1 in order of appearance, whereas the article prints the same items with the numbering of its own full document.  The complete native source of the article as distributed in the arXiv e-print for this exact version is bundled unchanged as \texttt{sources/152228/source0.tex}.
\begin{align}\label{eq 1.10}
    \Phi_{\mu*\mu}(z) := \sum_{j=1}^\infty \sum_{n=1}^\infty \frac{(\mu*\mu)(n)}{j} z^{jn}, \quad \lvert z \rvert <1.
\end{align}

\begin{align}\label{Major and Minor Arcs}
    \mathfrak{M} = \bigcup_{\substack{1 \leqslant a \leqslant q \leqslant Q \\ (a,q)=1}} \mathfrak{M}(q,a) \quad \textnormal{and} \quad \mathfrak{m} =  \mathcal{U} \backslash \mathfrak{M}.
\end{align}

\begin{lemma}\label{Partial Summation Lemma}
Let $X \geqslant 2$ and suppose $B>0$ is fixed. Let $\theta \in \mathbb{R}$ such that $\theta = a/q+\beta$ where $a,q \in \mathbb{Z}$ with $0 \leqslant a < q \leqslant(\log X)^B$, $(a,q)=1$ and $\lvert \beta \rvert \leqslant X ^{-1}(\log X)^B$. Then there exists a constant $c_B>0$ depending only on $B$ such that
\begin{align*}
    \sum_{n\leqslant X }(\mu*\mu)(n)\e(n\theta) \ll_B X\exp(-c_B\sqrt{\log X}).
\end{align*}
\end{lemma}

\begin{lemma}\label{Major Arcs Theorem}
Let $\theta \in \mathfrak{M}$ as defined in \eqref{Major and Minor Arcs}. Then for any fixed $A>0$,
\begin{align*}
        \Phi_{\mu * \mu}(\rho \e (\theta)) \ll_A \frac{X}{(\log X)^{A}}.
\end{align*}
\end{lemma}

\begin{proof} Since $\rho = \exp(-1/X)$, it follows from \eqref{eq 1.10} that
\begin{align*}
    \Phi_{\mu * \mu}(\rho \e(\theta)) = \sum_{j=1}^\infty\sum_{n=1}^\infty \frac{1}{j} (\mu * \mu)(n) e^{-nj / X} \e(j n \theta).
\end{align*}
We replace the exponentially damped term above by the integral
\begin{align*}
    e^{-nj / X} = \int_n^\infty jX^{-1}e^{-tj/X}\, dt,
\end{align*}
leaving us with
\begin{align*}
    \Phi_{\mu * \mu}(\rho \e(\theta)) = \sum_{j=1}^\infty \frac{1}{j} \int_1^\infty jX^{-1}e^{-tj/X} \sum_{n \leqslant t}  (\mu * \mu)(n) \e(j n \theta) \, dt.
\end{align*}
Suppose $J=(\log X)^{2A}$. Trivially, we obtain
\begin{align}\label{Integral Truncating}
\sum_{j>J}\frac{1}{j} \int_1^\infty jX^{-1}e^{-tj/X} &\sum_{n \leqslant t}  (\mu * \mu)(n) \e(j n \theta) \, dt \notag \\
   & \ll \sum_{j>J}^\infty \frac{1}{j} \int_1^\infty jX^{-1}e^{-tj/X} t \log t \, dt  \notag \\
   &\ll \sum_{j>J}^\infty \frac{X \log X}{j^2} \int_{j/X}^\infty ue^{-u} \log u  \, du \ll_A \frac{X}{(\log X)^{A}}.
\end{align}
Hence, in what follows, it suffices to work with the sum
\begin{align}\label{Truncated J}
\sum_{j=1}^J \frac{1}{j} \int_1^\infty jX^{-1}e^{-tj/X} \sum_{n \leqslant t}  (\mu * \mu)(n) \e(j n \theta) \, dt.
\end{align}
Writing $\theta=a/q+\beta$ with $q\leqslant(\log X)^A$ and $|\beta|\leqslant\delta_q \leqslant X ^{-1}(\log X)^A$, one has
\[ \theta _j := j\theta = \frac{aj/(j,q)}{q/(j,q)}+ j\beta = \frac{a _j}{q _j}+\beta _j,
\]
where $(a _j,q _j)=1$, $1 \leqslant q _j \leqslant q \leqslant(\log X)^{A} $ and $\beta _j \leqslant X ^{-1}(\log X)^{3A}$. We can now break the inner integral in \eqref{Truncated J} as
\begin{align}
  I &= \int_1^\infty jX^{-1}e^{-tj/X} \sum_{n \leqslant t}  (\mu * \mu)(n) \e(j n \theta) \, dt \notag \\
  &= \bigg(\int_1^{X \log X} + \int_{X\log X}^\infty\bigg) jX^{-1}e^{-tj/X} \sum_{n \leqslant t}  (\mu * \mu)(n) \e(j n \theta) \, dt = I_1+I_2 \notag, 
\end{align}
say. For $I_2$, one has
\begin{align}
    I_2 & \ll   \int_{X\log X}^\infty jX^{-1}e^{-tj/X} t \log t \, dt  \ll \frac{X \log X}{j} \int_{\log X}^\infty ue^{-u} \log u  \, du \ll_A \frac{X}{(\log X)^{3A}}. \label{I_2 Estimate}
\end{align}
For $I_1$, an application of Lemma \ref{Partial Summation Lemma} shows that
\begin{align}\label{I_1 Estimate}
    I_1 & \ll_A X \log X \exp(-C_A \sqrt{\log X})  \int_1^{X\log X} jX^{-1}e^{-tj/X} \, dt \notag \\
   & \ll_A X \log X \exp(-C_A \sqrt{\log X}) \int_{1}^{X\log X} e^{-u}   \, du \ll_A \frac{X}{(\log X)^{3A}},
\end{align}
where $C_A>0$ is a constant depending only on $A$. Putting together \eqref{I_2 Estimate} and \eqref{I_1 Estimate}, we obtain
\[ \sum_{j=1}^J \frac{1}{j} \int_1^\infty jX^{-1}e^{-tj/X} \sum_{n \leqslant t}  (\mu * \mu)(n) \e(j n \theta) \, dt \ll_A \frac{X}{(\log X)^{A}}.
\]
Combining this with \eqref{Integral Truncating}, we arrive at our desired conclusion.
\end{proof}

\end{document}
